\medskip

\textbf{PARTIE A}

\medskip

On définit sur l'intervalle $]0~;~+\infty[$ la fonction $g$ par : 
\[g(x) = \dfrac 2x - \dfrac{1}{x^2} + \ln x \quad \text{où ln désigne la fonction logarithme népérien}.\]

On admet que la fonction $g$ est dérivable sur $]0~;~+\infty[ = I$ et on note $g'$ sa fonction dérivée.

\medskip

\begin{enumerate}
\item Montrer que pour $x > 0$, le signe de $g'(x)$ est celui du trinôme du second degré $\left(x^2 - 2x + 2\right)$.
\item En déduire que la fonction $g$ est strictement croissante sur $]0~;~+\infty[$.
\item Montrer que l'équation $g(x) = 0$ admet une unique solution sur l'intervalle [0,5~;~1], que l'on notera $\alpha$.
\item On donne le tableau de signes de $g$ sur l'intervalle $]0~;~+\infty[ = I$ :
\begin{center}
\psset{unit=1cm}
\begin{pspicture}(4,1)
\psline(0,0.5)(4,0.5)\psline(1,0)(1,1)\psline(1.1,0)(1.1,0.5)
\uput[u](0.5,0.4){$x$} \uput[u](1.2,0.4){$0$} \uput[u](2.5,0.4){$\alpha$} \uput[u](3.7,0.4){$+\infty$}
\uput[u](0.5,-0.1){$g(x)$}\uput[u](1.75,-0.1){$-$}\uput[u](2.5,-0.1){$0$}
\uput[u](3.25,-0.1){$+$}
\end{pspicture}
\end{center}

Justifier ce tableau de signes à l'aide des résultats obtenus aux questions précédentes.
\end{enumerate}

\medskip

\textbf{PARTIE B}

\medskip

On considère la fonction $f$ définie sur l'intervalle $]0~;~+\infty[ = I$ par: 
\[f(x) = \e^x \ln x.\]

On note $\mathcal{C}_f$ la courbe représentative de $f$ dans un repère orthonormé.

\medskip

\begin{enumerate}
\item On admet que la fonction $f$ est deux fois dérivable sur $]0~;~+\infty[$, on note $f'$ sa fonction dérivée, $f''$ sa fonction dérivée seconde et on admet que :

pour tout nombre réel $x > 0,\:\: f'(x) = \e^x \left(\dfrac 1x + \ln x\right)$.

Démontrer que, pour tout nombre réel $x > 0$, on a : $f''(x) = \e^x \left(\dfrac 2x - \dfrac{1}{x^2} + \ln x\right)$.
\item On pourra remarquer que pour tout réel $x > 0$,\: $f''(x) = \e^x \times  g(x)$, où $g$ désigne la fonction étudiée dans la partie A.
\item 
	\begin{enumerate}
		\item Dresser le tableau de signes de la fonction $f''$ sur $]0~;~+\infty[$. Justifier.
		\item Justifier que la courbe $\mathcal{C}_f$ admet un unique point d'inflexion A.
		\item Étudier la convexité de la fonction $f$ sur l'intervalle $]0~;~+\infty[$. Justifier.
	\end{enumerate}
\item
	\begin{enumerate}
		\item Calculer les limites de $f$ aux bornes de son ensemble de définition.
		\item Montrer que $f'(\alpha) = \dfrac{\e^{\alpha}}{\alpha^2}(1 - \alpha)$.

On rappelle que $\alpha$ est l'unique solution de l'équation $g(x) = 0$.
		\item Démontrer que $f'(\alpha) > 0$ et en déduire le signe de $f'(x)$ pour $x$ appartenant à $]0~;~+\infty[$.
		\item En déduire le tableau de variations complet de la fonction $f$ sur $]0~;~+\infty[$.
	\end{enumerate}
\end{enumerate}
